% !TEX program = tectonic
\documentclass[11pt,a4paper]{article}
\newcommand{\DocumentKind}{Documentação técnica}
\newcommand{\DocumentVersion}{1.4}
\usepackage{gemdados}

\begin{document}

\GemCover
  {Documentação Técnica}
  {Plataforma de dados, infraestrutura e operação}
  {Descrição consolidada da infraestrutura, pipelines, views analíticas e dashboard Looker Studio.}

\section*{Controle do documento}

\begin{tabularx}{\textwidth}{L{35mm}Y}
\toprule
\textbf{Campo} & \textbf{Valor} \\
\midrule
Versão & 1.4 \\
Data de referência & 24 de setembro de 2026 \\
Público & Time técnico, mantenedores, mentores e novos colaboradores \\
Escopo & Infraestrutura GCP, pipelines, views analíticas no BigQuery e dashboard Looker Studio \\
\bottomrule
\end{tabularx}

\subsection*{Critério}

O conteúdo descreve os recursos provisionados em nuvem, repositórios de código,
views analíticas SQL no BigQuery e a integração com a camada de visualização
no Looker Studio em 24 de setembro de 2026.

\clearpage
\tableofcontents
\clearpage

\section{Objetivo e escopo}

A plataforma de dados em nuvem GEM Dados cobre infraestrutura como código, serviço
de ingestão conteinerizado, modelagem analítica no BigQuery e camada de visualização
no Looker Studio.

Este documento registra:

\begin{itemize}
  \item a responsabilidade de cada repositório e camada;
  \item os recursos de infraestrutura definidos para os ambientes;
  \item as rotas, funções e variáveis do serviço de ingestão;
  \item as declarações de fontes e tabelas staging;
  \item as views analíticas SQL que centralizam as regras de cálculo;
  \item a arquitetura de consumo do dashboard Looker Studio;
  \item as esteiras de validação, deploy e controles de segurança.
\end{itemize}

\section{Termos essenciais}

\begin{tabularx}{\textwidth}{L{32mm}Y}
\toprule
\textbf{Termo} & \textbf{Significado neste projeto} \\
\midrule
GCP & Google Cloud Platform, conjunto de serviços em nuvem usado pela solução. \\
GCS & Google Cloud Storage, armazenamento de objetos usado para a landing zone. \\
BigQuery & Banco analítico que hospeda os datasets \code{raw}, \code{staging} e \code{marts}. \\
Cloud Run & Serviço que executa o container Python de ingestão. \\
Dataform & Serviço que compila e executa as definições SQLX no BigQuery. \\
Terraform & Ferramenta de infraestrutura como código usada em \code{cloud\_iac}. \\
Cloud Build & Serviço de CI/CD que valida, constrói e publica mudanças. \\
IAM & Controle de identidades e permissões dos recursos do GCP. \\
Looker Studio & Ferramenta de BI que provê o dashboard de acompanhamento de mentorados. \\
Views analíticas & Camada SQL no BigQuery que centraliza regras e métricas calculadas. \\
\bottomrule
\end{tabularx}

\section{Arquitetura registrada no código}

\subsection{Componentes}

\begin{enumerate}
  \item O repositório \code{cloud\_iac} define recursos e permissões no GCP com
  Terraform.
  \item O repositório \code{data\_ingestion} contém uma aplicação Flask
  configurada para Cloud Run.
  \item O repositório \code{data\_models} registra dez fontes raw no Dataform.
  \item Cloud Build, Cloud Workflows e Cloud Scheduler conectam deploy e
  execução conforme os arquivos Terraform e YAML versionados.
  \item Views SQL analíticas intermediárias consolidam as regras de negócio
  sobre a camada \code{staging}.
  \item O Looker Studio consome as views e tabelas de staging para alimentar
  o Dashboard de Acompanhamento de Mentorados.
\end{enumerate}

\subsection{Fluxo da arquitetura de ponta a ponta}

\begin{center}
\code{Staging (BigQuery) -> Views SQL Analíticas -> Looker Studio -> Dashboard}
\end{center}

\subsection{Ambientes}

\begin{tabularx}{\textwidth}{L{25mm}L{46mm}L{25mm}Y}
\toprule
\textbf{Ambiente} & \textbf{Projeto GCP} & \textbf{Branch} & \textbf{Finalidade registrada} \\
\midrule
Staging & \code{gem-dados-lake-stg} & \code{stg} & Ambiente de validação integrado e origem do dashboard. \\
Produção & \code{gem-dados-lake-prd} & \code{main} & Ambiente associado à branch principal. \\
\bottomrule
\end{tabularx}

Cada ambiente possui diretório Terraform, backend de state, projeto GCP e
valores próprios. A promoção descrita nos READMEs ocorre de \code{stg} para
\code{main} por pull request.

\section{Mapa dos repositórios}

\begin{longtable}{L{34mm}L{35mm}L{47mm}L{33mm}}
\toprule
\textbf{Repositório} & \textbf{Tecnologia} & \textbf{Responsabilidade} & \textbf{Artefato produzido} \\
\midrule
\endhead
\code{cloud\_iac} & Terraform e Cloud Build & Define recursos, IAM, conexões e esteiras no GCP. & Plano e aplicação de infraestrutura. \\
\code{data\_ingestion} & Python, Flask e Cloud Run & Expõe rotas HTTP e executa jobs de ingestão. & Imagem Docker para Cloud Run. \\
\code{data\_models} & SQLX e Dataform & Declara fontes e transforma dados no BigQuery. & Views de staging e tabelas de marts. \\
\bottomrule
\end{longtable}

\section{Infraestrutura - \code{cloud\_iac}}

\subsection{Diretórios de configuração}

\begin{tabularx}{\textwidth}{L{34mm}Y}
\toprule
\textbf{Diretório} & \textbf{Conteúdo versionado} \\
\midrule
\code{bootstrap/} & Bucket de state, conta \code{terraform-ci}, Secret Manager, conexão GitHub 2nd gen e triggers iniciais. \\
\code{org/} & Tag de governança \code{environment}, valores de ambiente e permissão \code{tagUser}. \\
\code{envs/stg/} & Composição Terraform do projeto \code{gem-dados-lake-stg}. \\
\code{envs/prd/} & Composição Terraform do projeto \code{gem-dados-lake-prd}. \\
\code{modules/} & Módulos Terraform chamados pelos ambientes e pelo bootstrap. \\
\bottomrule
\end{tabularx}

\clearpage
\subsection{Módulos Terraform presentes no repositório}

\begin{longtable}{L{43mm}L{48mm}L{58mm}}
\toprule
\textbf{Módulo} & \textbf{Recurso principal} & \textbf{Configuração encontrada} \\
\midrule
\endhead
\code{project\_baseline} & APIs e tag de ambiente & Habilita APIs e pode vincular a tag de governança ao projeto. \\
\code{gcs\_bucket} & Bucket GCS & Acesso uniforme, prevenção de acesso público, versionamento e lifecycle. \\
\code{bigquery\_dataset} & Dataset BigQuery & Localização, descrição, labels e controles de destruição. \\
\code{artifact\_registry} & Repositório Docker & Repositório usado para a imagem de ingestão. \\
\code{cloud\_run\_service} & Serviço Cloud Run v2 & Imagem, conta de serviço, recursos, env vars, secrets e invocadores. \\
\code{dataform\_repository} & Repositório Dataform & URL Git, branch, token e banco padrão. \\
\code{dataform\_orchestration} & Workflow e Scheduler & Conta orquestradora, permissões, Workflow e agendamento. \\
\code{ingestion\_orchestration} & Workflow e Scheduler & Orquestração da ingestão conteinerizada antes do Dataform. \\
\code{cloudbuild\_connection} & Conexão GitHub 2nd gen & Conexão e registro de repositório no Cloud Build. \\
\code{cloudbuild\_trigger} & Trigger Cloud Build & Evento push ou pull request, branch, arquivo YAML e substituições. \\
\bottomrule
\end{longtable}

\subsection{Recursos definidos por ambiente}

Os arquivos \code{envs/stg/main.tf} e \code{envs/prd/main.tf} possuem a mesma
estrutura. Neles estão definidos:

\begin{itemize}
  \item habilitação de APIs e vínculo da tag de ambiente;
  \item conta de serviço \code{data-ingestion} e seus vínculos IAM;
  \item datasets BigQuery \code{raw}, \code{staging} e \code{marts};
  \item bucket \code{<project>-raw};
  \item Artifact Registry \code{data-ingestion};
  \item serviço Cloud Run \code{data-ingestion};
  \item repositório Dataform \code{data-models};
  \item conta \code{dataform-runner} e permissões de BigQuery;
  \item Workflow e Scheduler do Dataform;
  \item registros e triggers Cloud Build de \code{data\_ingestion} e
  \code{data\_models};
  \item permissão para a esteira invocar o Workflow.
\end{itemize}

\subsection{Valores dos ambientes}

\begin{tabularx}{\textwidth}{L{34mm}L{47mm}Y}
\toprule
\textbf{Variável} & \textbf{Valor versionado} & \textbf{Uso} \\
\midrule
\code{project\_id} & \shortstack[l]{\code{gem-dados-lake-stg}\\\code{gem-dados-lake-prd}} & Projeto do ambiente. \\
\code{env} & \code{stg} ou \code{prd} & Nomes, labels e branch de deploy. \\
\code{region} & \code{us-central1} & Região dos serviços regionais. \\
\code{bq\_location} & \code{US} & Localização dos datasets BigQuery. \\
\code{github\_owner} & \code{gem-dados} & Organização dos repositórios. \\
\code{org\_id} & \code{809669352691} & Organização usada pela tag de governança. \\
\bottomrule
\end{tabularx}

\subsection{Esteiras de infraestrutura}

\subsubsection{Pull request}

O arquivo \code{cloudbuild-pr.yaml} executa \code{gitleaks}, \code{tfsec},
\code{terraform fmt -check}, \code{terraform validate} e \code{terraform plan}.
O YAML não possui etapa de apply.

\subsubsection{Push na branch do ambiente}

O arquivo \code{cloudbuild.yaml} executa gitleaks, validação, plan e
\code{terraform apply}. A variável \code{\_ENV} seleciona \code{envs/stg} ou
\code{envs/prd}.

\section{Ingestão - \code{data\_ingestion}}

\subsection{Aplicação HTTP}

O arquivo \code{src/main.py} cria uma aplicação Flask com um registro interno de
jobs. As rotas existentes são:

\begin{tabularx}{\textwidth}{L{38mm}L{24mm}Y}
\toprule
\textbf{Rota} & \textbf{Método} & \textbf{Comportamento no código} \\
\midrule
\code{/healthz} & GET & Retorna um JSON com status \code{ok}. \\
\code{/run/<job\_name>} & POST & Procura o job no dicionário \code{JOBS}, executa a rotina correspondente e retorna \code{done}; nomes ausentes retornam HTTP 404. \\
\bottomrule
\end{tabularx}

\subsection{Jobs de ingestão registrados}

\begin{longtable}{L{48mm}L{100mm}}
\toprule
\textbf{Nome no endpoint} & \textbf{Script e função executada} \\
\midrule
\endhead
\code{ingest} & \code{identify\_datacamp\_csvs.py}: identifica dinamicamente os relatórios no Drive por assinatura de colunas. \\
\code{upload\_gcs} & \code{upload\_parquet\_to\_gcs.py}: converte os relatórios para Parquet e realiza o upload para o GCS. \\
\code{load\_bq} & \code{load\_gcs\_parquet\_to\_bigquery.py}: aterrissa os dados Parquet do GCS nas tabelas \code{raw} do BigQuery. \\
\code{rotina\_sem\_duplicatas} & \code{rotina\_sem\_duplicatas\_job.py}: pipeline integrado e idempotente sem geração de registros duplicados. \\
\bottomrule
\end{longtable}

O Dockerfile usa Python 3.12 slim, instala as dependências, copia \code{src/},
cria o usuário não-root \code{appuser} e inicia Gunicorn na porta definida por
\code{PORT}.

\subsection{Variáveis lidas pela aplicação}

\begin{tabularx}{\textwidth}{L{34mm}L{26mm}Y}
\toprule
\textbf{Variável} & \textbf{Leitura} & \textbf{Uso no código} \\
\midrule
\code{GCP\_PROJECT} & Obrigatória & Projeto dos clientes GCS e BigQuery. \\
\code{RAW\_BUCKET} & Obrigatória & Bucket da landing zone. \\
\code{ENVIRONMENT} & Padrão \code{stg} & Campo \code{env} do payload. \\
\code{BQ\_DATASET} & Padrão \code{raw} & Dataset da tabela carregada. \\
\code{PORT} & Padrão \code{8080} & Porta HTTP do servidor. \\
\bottomrule
\end{tabularx}

\subsection{Esteira de deploy}

O \code{cloudbuild.yaml} executa gitleaks, constrói a imagem Docker, publica a
tag \code{SHORT\_SHA} no Artifact Registry e atualiza o serviço Cloud Run com
\code{--no-allow-unauthenticated}.

\section{Modelagem e Views Analíticas}

\subsection{Configuração Dataform}

O arquivo \code{workflow\_settings.yaml} registra:

\begin{itemize}
  \item projeto padrão \code{gem-dados-lake-stg};
  \item localização BigQuery \code{US};
  \item dataset padrão \code{staging};
  \item dataset de assertions \code{dataform\_assertions};
  \item Dataform Core 3.0.0;
  \item variável \code{env} com valor \code{stg}.
\end{itemize}

\subsection{Declarações de fontes}

\begin{longtable}{L{48mm}L{57mm}L{44mm}}
\toprule
\textbf{Nome Dataform} & \textbf{Origem citada na descrição} & \textbf{Descrição registrada} \\
\midrule
\endhead
\code{raw\_members} & \code{membros\_export} & Membros e licenças. \\
\code{raw\_summary} & \code{resumo\_export} & Resumo geral da plataforma. \\
\code{raw\_certifications} & \shortstack[l]{\code{resumo\_por\_certificacao}\\\code{\_export}} & Progresso por certificação. \\
\code{raw\_content\_catalog} & \code{catalogo\_de\_conteudo\_export} & Catálogo de conteúdos. \\
\code{raw\_course\_activity} & \code{resumo\_por\_curso\_export} & Atividade por curso. \\
\code{raw\_program\_activity} & \code{resumo\_por\_programa\_export} & Atividade por programa ou trilha. \\
\code{raw\_project\_activity} & \code{resumo\_por\_projeto\_export} & Atividade por projeto prático. \\
\code{raw\_skill\_assessments} & \shortstack[l]{\code{avaliacao\_de\_habilidades}\\\code{\_export}} & Avaliações de habilidades. \\
\code{raw\_team\_history} & \code{historico\_da\_equipe\_export} & Histórico das equipes. \\
\code{raw\_time\_spent} & \code{tempo\_no\_aprendizado\_export} & Tempo de aprendizado. \\
\bottomrule
\end{longtable}

Os dez arquivos iniciados por \code{raw\_} são declarações Dataform do tipo
\code{declaration}.

\subsection{Camada de staging (\code{definitions/staging})}

A camada staging padroniza nomes de colunas, aplica \code{SAFE\_CAST} para
conversão estrita de tipos, normaliza nulos e substitui dados identificáveis
pelo código anônimo \code{user\_id}. Conta com dez modelos:

\begin{itemize}
  \item \code{stg\_members}, \code{stg\_summary}, \code{stg\_certifications}, \code{stg\_content\_catalog};
  \item \code{stg\_course\_activity}, \code{stg\_program\_activity}, \code{stg\_project\_activity};
  \item \code{stg\_skill\_assessments}, \code{stg\_team\_history}, \code{stg\_time\_spent}.
\end{itemize}

\subsection{Camada de marts (\code{definitions/marts})}

A camada de marts organiza tabelas fato e dimensão para consumo analítico:

\begin{itemize}
  \item \textbf{Dimensões}: \code{dim\_usuarios}, \code{dim\_conteudo};
  \item \textbf{Fatos}: \code{fct\_progresso\_cursos}, \code{fct\_progresso\_trilhas}, \code{fct\_atividades\_projetos}, \code{fct\_certificacoes}, \code{fct\_avaliacoes\_habilidades}, \code{fct\_resumo\_usuario}, \code{fct\_historico\_equipe}, \code{fct\_tempo\_aprendizado};
  \item \textbf{Agregações periódicas}: \code{mart\_metricas\_semanais}.
\end{itemize}

\subsection{Testes de qualidade e integridade}

São executadas assertions automáticas no Dataform para garantir consistência:
\begin{itemize}
  \item \code{assert\_raw\_stg\_row\_count}: validação de contagem de linhas entre raw e staging;
  \item \code{assert\_raw\_stg\_metrics}: validação de consistência e completude métrica.
\end{itemize}

\subsection{Views analíticas no BigQuery (\code{gem-dados-lake-stg.staging})}

Para centralizar as regras de negócio e desacoplar o cálculo das ferramentas de
visualização, foram criadas views analíticas nativas no BigQuery:

\begin{longtable}{L{45mm}L{45mm}L{58mm}}
\toprule
\textbf{View analítica} & \textbf{Tabela base} & \textbf{Indicadores consolidados} \\
\midrule
\endhead
\code{vw\_program\_activity} & \code{stg\_program\_activity} & Taxas de conclusão de assessments, capítulos, cursos e projetos práticos. \\
\bottomrule
\end{longtable}

\subsection{Decisão técnica: migração dos cálculos para SQL}

Inicialmente, os cálculos de taxas e proporções estavam sendo implementados na
camada semântica do Looker Studio. Seguindo recomendação do Product Owner (PO),
essas regras foram migradas para views SQL nativas no BigQuery.

Essa abordagem teve como benefícios:
\begin{itemize}
  \item centralização e governança das regras de negócio;
  \item eliminação do acoplamento com a ferramenta de visualização;
  \item possibilidade de reaproveitamento das métricas por outras ferramentas de BI;
  \item rastreabilidade direta do cálculo até a camada de staging;
  \item simplificação da camada de apresentação no Looker Studio.
\end{itemize}

\section{Camada de visualização - Looker Studio}

\subsection{Conexões e origens do dashboard}

O Dashboard de Acompanhamento de Mentorados está hospedado no Looker Studio e
conecta diretamente ao dataset \code{staging} do projeto \code{gem-dados-lake-stg}:

\begin{itemize}
  \item \code{vw\_program\_activity}: taxas de conclusão de atividades (cursos, capítulos, projetos e assessments);
  \item \code{stg\_certifications}: certificações obtidas e status;
  \item \code{stg\_skill\_assessment}: avaliações e proficiência técnica;
  \item \code{stg\_project\_activity}: conclusões de projetos práticos.
\end{itemize}

\subsection{Estrutura de páginas e controles}

O painel é estruturado em três páginas:
\begin{enumerate}
  \item \textbf{Progresso \& Volume}: cards consolidados de conclusões (trilhas, alunos, cursos, projetos, assessments), distribuição por tipo de atividade e taxas de conclusão;
  \item \textbf{Qualificação}: certificações totais, assessments realizados, desempenho por habilidade e entregas agrupadas por equipe;
  \item \textbf{Desempenho}: médias de cursos por trilha e taxa de conclusão de cursos por trilha.
\end{enumerate}

Controles dinâmicos de filtro foram aplicados para Trilha (\textit{Track}), Período (\textit{Period}) e Nível de Habilidade (\textit{Skill Level}).

\section{Operação registrada}

\subsection{Fluxo Git descrito nos READMEs}

\begin{lstlisting}
git checkout stg
git pull
git checkout -b feat/minha-mudanca
pre-commit run --all-files
git push -u origin feat/minha-mudanca
\end{lstlisting}

Os READMEs descrevem pull request para \code{stg}, validação no ambiente e
promoção posterior para \code{main}.

\subsection{Pontos de observação dos serviços definidos}

\begin{tabularx}{\textwidth}{L{36mm}L{50mm}Y}
\toprule
\textbf{Componente} & \textbf{Local} & \textbf{Registro gerado pelo código} \\
\midrule
Cloud Build & Builds do projeto GCP & Etapas de validação, build, push, deploy ou Terraform. \\
Cloud Run & Logs do serviço & Logs Python do job e respostas HTTP. \\
Workflows & Workflow \code{dataform-<env>} & Compilação e criação da invocação Dataform. \\
BigQuery & Jobs e Views & Execução de consultas analíticas e atualização de views. \\
Looker Studio & Relatório de BI & Métricas visuais e logs de conexão com BigQuery. \\
\bottomrule
\end{tabularx}

\section{Controles de segurança versionados}

\begin{itemize}
  \item gitleaks nos hooks de pre-commit e nos arquivos Cloud Build;
  \item tfsec, terraform fmt, validate e plan na esteira de pull request;
  \item \code{detect-private-key} nos hooks de infraestrutura e ingestão;
  \item arquivos de credencial, state e chaves bloqueados pelos \code{.gitignore};
  \item Secret Manager para o token da conexão GitHub;
  \item contas de serviço separadas para Terraform, ingestão, orquestração e
  execução Dataform;
  \item bucket raw com prevenção de acesso público;
  \item Cloud Run implantado sem acesso anônimo;
  \item container de ingestão executado por usuário não-root;
  \item anonimização completa dos usuários com uso de \code{user\_id} em todas as views e telas de BI.
\end{itemize}

\section{Inventário dos componentes}

\begin{tabularx}{\textwidth}{L{48mm}Y}
\toprule
\textbf{Área} & \textbf{Conteúdo presente} \\
\midrule
Infraestrutura & Bootstrap, configuração de organização, dois ambientes, dez módulos Terraform (incluindo orquestração) e esteiras de plan/apply. \\
Ingestão & Aplicação Flask/Cloud Run, quatro jobs registrados em \code{src/jobs/}, Dockerfile e esteira de deploy Cloud Run. \\
Modelagem & Dez declarações raw, modelos de staging (\code{stg\_*}), marts (\code{dim\_*}, \code{fct\_*}), asserções de qualidade, configuração Dataform e esteira de execução. \\
Views analíticas & View SQL \code{vw\_program\_activity} no BigQuery staging. \\
Visualização & Dashboard Looker Studio com três páginas (Progresso, Qualificação, Desempenho) e 3 filtros. \\
Segurança & Anonimização por \code{user\_id}, gitleaks, tfsec, Bandit/Ruff, hooks pre-commit e IAM em Terraform. \\
\bottomrule
\end{tabularx}

\section{Localização dos componentes}

\begin{tabularx}{\textwidth}{L{46mm}Y}
\toprule
\textbf{Componente} & \textbf{Arquivo, serviço ou diretório} \\
\midrule
Composição dos ambientes & \code{cloud\_iac/envs/stg} e \code{cloud\_iac/envs/prd} \\
Módulos Terraform & \code{cloud\_iac/modules} (10 módulos) \\
Servidor e jobs de ingestão & \code{data\_ingestion/src/main.py} e \code{data\_ingestion/src/jobs/} \\
Imagem do serviço & \code{data\_ingestion/Dockerfile} \\
Declarações raw & \code{data\_models/definitions/sources} \\
Modelos Dataform & \code{data\_models/definitions/staging} e \code{data\_models/definitions/marts} \\
Views analíticas SQL & BigQuery: \code{gem-dados-lake-stg.staging.vw\_program\_activity} \\
Camada de BI & Looker Studio: Dashboard de Acompanhamento de Mentorados \\
\bottomrule
\end{tabularx}

\end{document}
